% ======================================================================
% preprints/jphq_06_structural_tension_nonoptimizable/sections/08_literature_positioning.tex
% Journal-grade content (100/100). ASCII-only.
% ======================================================================

\section{Positioning with Respect to Existing Literature}

This section situates the present result with respect to adjacent bodies
of work. Its purpose is strictly delimiting rather than comparative or
evaluative. We identify classes of approaches that are formally
incompatible with the notion of structural tension as defined in OC/ETS
and clarify the source of that incompatibility.

No claim of priority, coverage, or superiority over existing literature
is made beyond the stated impossibility results.

\subsection{Optimization and Control Frameworks}

Classical optimization and control frameworks presuppose the existence of
an objective function, an ordered state space, or both. Structural
constraints are typically represented as penalties, soft constraints, or
barrier terms embedded in a metric setting.

In OC/ETS, thresholds play a categorically different role. They are hard
admissibility conditions that delimit whether the enterprise continuum
$\KEA$ is defined at all. Thresholds do not degrade performance; they
invalidate the object of diagnosis.

Consequently, optimization-based readings of structural tension are not
approximations or relaxations of the OC/ETS formalism. They instantiate a
different formal object with different semantics.

\subsection{Resilience and Systems Engineering}

In resilience-oriented literature, resilience is commonly defined as the
capacity to absorb disturbance and return to a reference state. Such
definitions presuppose both a metric of deviation and a privileged
baseline, attractor, or equilibrium region.

Structural tension assumes neither. It encodes no notion of distance,
deviation, recovery, or return. It evaluates only whether admissible
diagnosis remains defined under the current structural constraints.

Accordingly, structural tension cannot be reduced to resilience indices,
robustness margins, or stability measures without introducing metric
structure absent from OC/ETS.

\subsection{Complexity and Phase-Transition Analogies}

Phase-transition metaphors are common in complexity science. In such
contexts, phases are typically characterized by continuous order
parameters, bifurcations, or critical behavior.

The phase labels used here (\PHASESUCCESS, \PHASEINERTIA,
\PHASECOLLAPSE, \PHASESTOP) are logical classes induced by threshold
inequalities and explicit precedence rules. They are not associated with
continuous parameters, critical exponents, or scaling laws.

Any analogy to physical phase transitions is therefore structural rather
than quantitative and must not be interpreted as implying continuity,
metric proximity, or dynamical bifurcation behavior.

\subsection{Maturity Models and Capability Scales}

Enterprise maturity models impose ordinal or cardinal scales that rank
organizations according to assumed progress, sophistication, or
capability accumulation.

The OC/ETS formalism explicitly rejects such ordering. Structural tension
does not rank enterprises, assess advancement, or imply improvement
paths. It determines only whether statements of a given diagnostic class
remain formally admissible.

Maturity levels, capability scores, and similar constructs are therefore
out of scope by construction.

\subsection{Formal Ontologies and Meta-Models}

Many enterprise ontologies and meta-models aim at semantic completeness or
exhaustive representational coverage. OC/ETS adopts the opposite stance:
it enforces minimality.

Only those primitives required to delimit admissibility and
diagnosability are introduced. No attempt is made to model enterprise
reality in full.

This minimality is essential to the impossibility results established in
this paper. Increasing representational richness typically reintroduces
metric, ordering, or optimization assumptions that OC/ETS excludes by
design.

\subsection{Summary of Structural Distinctions}

The structural-tension formalism differs from adjacent approaches along
three non-negotiable dimensions:
(i) thresholds are admissibility boundaries, not costs or penalties;
(ii) phases are logical classes, not ordered levels;
(iii) no scalar, metric, or optimization interpretation is defined.

These distinctions are formal rather than rhetorical. Violating any of
them results in a different theory object, not an alternative reading of
OC/ETS.

%%% End of section %%%
